\documentclass[11pt,oneside]{book}

\usepackage[a4paper,margin=1in]{geometry}
\usepackage{amsmath,amssymb,amsthm,mathtools}
\usepackage{enumitem}
\usepackage{array,longtable,booktabs}
\usepackage{xcolor}
\usepackage{hyperref}

\hypersetup{
  colorlinks=true,
  linkcolor=blue!55!black,
  urlcolor=blue!55!black,
  citecolor=blue!55!black
}
\emergencystretch=2em

\newtheorem{theorem}{Theorem}[chapter]
\newtheorem{proposition}[theorem]{Proposition}
\newtheorem{lemma}[theorem]{Lemma}
\newtheorem{corollary}[theorem]{Corollary}
\newtheorem{definition}[theorem]{Definition}
\newtheorem{remark}[theorem]{Remark}
\newtheorem{example}[theorem]{Example}

\DeclareMathOperator{\Tr}{Tr}
\DeclareMathOperator{\diag}{diag}
\DeclareMathOperator{\Var}{Var}
\DeclareMathOperator{\dist}{dist}
\DeclareMathOperator{\Lip}{Lip}
\DeclareMathOperator{\card}{card}
\DeclareMathOperator{\med}{med}
\DeclareMathOperator{\id}{id}
\DeclareMathOperator{\rank}{rank}
\DeclareMathOperator{\supp}{supp}
\DeclareMathOperator{\Mass}{Mass}
\DeclareMathOperator{\Haar}{Haar}
\DeclareMathOperator{\vol}{vol}
\DeclareMathOperator{\End}{End}
\DeclareMathOperator{\Hom}{Hom}

\newcommand{\C}{\mathbb C}
\newcommand{\R}{\mathbb R}
\newcommand{\N}{\mathbb N}
\newcommand{\E}{\mathbb E}
\newcommand{\Prob}{\mathbb P}
\newcommand{\1}{\mathbf 1}
\newcommand{\Sphere}{\mathbb S}
\newcommand{\eps}{\varepsilon}
\newcommand{\pt}{\Gamma}
\newcommand{\HS}{\mathrm{HS}}
\newcommand{\op}{\mathrm{op}}
\newcommand{\dd}{\,\mathrm d}
\newcommand{\abs}[1]{\left|#1\right|}
\newcommand{\norm}[1]{\left\|#1\right\|}
\newcommand{\ip}[2]{\left\langle #1,#2\right\rangle}
\newcommand{\set}[1]{\left\{#1\right\}}
\newcommand{\given}{\,:\,}
\newcommand{\leanfile}[1]{\begingroup\scriptsize\texttt{#1}\endgroup}

\title{Appendix B Formal Concordance Draft}
\author{An audit-oriented translation of the formalized Appendix B}
\date{\today}

\begin{document}
\frontmatter
\maketitle

\chapter*{Preface}

This document is now kept as an audit-oriented concordance draft for the
formalized Appendix B.  It records the mathematical blocks exposed by the
formal development and keeps them close to the structure of the proof scripts:
finite index bookkeeping, normalizations, measurability, null sets, medians,
radial--spherical factorizations, Wick pairings, defect counts, and the final
concentration assembly.

It is not the reader-first manuscript.  The polished mathematical exposition is
the separate file \texttt{AppendixB\_bookstyle.tex}.  The present file remains
useful as a bridge between the formal modules and the expository rewrite.  In
particular, the following three facts are treated as closed mathematical
theorems here because they are closed in the formal development:
\[
  \text{polar law},\qquad
  \text{radius--direction independence},\qquad
  \text{strong global surface Levy concentration}.
\]

\tableofcontents

\mainmatter

\chapter{The Finite Model}

\section{Finite index sets and matrix spaces}

Let \(P,Q,\Sigma\) be finite sets.  The intended balanced specialization is
\[
  P=Q=\{1,\ldots,d\},\qquad \Sigma=\{1,\ldots,s\},
\]
but most of the argument is carried out for arbitrary finite sets.  Write
\[
  I=P\times Q,\qquad D=\card I=\card P\,\card Q,\qquad S=\card\Sigma .
\]
Whenever an asymptotic estimate contains a parameter \(\lambda\), we are in
the balanced asymptotic regime
\[
  P=Q=\{1,\ldots,d\},\qquad \Sigma=\{1,\ldots,s_d\},\qquad
  \frac{s_d}{d^2}\longrightarrow \lambda\in(0,\infty).
\]
Thus \(D=d^2\) and \(S=s_d\).  Constants denoted \(C_\lambda\), \(c_\lambda\),
or similar may depend on \(\lambda\), but not on \(d\).

The Hilbert space of sample matrices is
\[
  \mathcal X=\C^{I\times \Sigma},
\]
which we identify with the space of linear maps
\(\C^\Sigma\to\C^I\).  Both \(\C^I\) and \(\C^\Sigma\) carry their standard
Hermitian inner products:
\[
  \ip{x}{y}=\sum_i x_i\overline{y_i}.
\]

\begin{definition}[Matrix norms]
Let \(E\) and \(F\) be finite-dimensional complex Hilbert spaces and let
\(T:E\to F\) be linear.  The operator norm of \(T\) is
\[
  \norm T_{\op}=\sup_{\norm v_2=1}\norm{Tv}_2.
\]
Equivalently, it is the largest singular value of \(T\).

The Hilbert--Schmidt norm is
\[
  \norm T_{\HS}^2=\Tr(T^*T),
\]
or, in orthonormal bases,
\[
  \norm T_{\HS}^2=\sum_{a,b}|T_{ab}|^2.
\]
The trace norm, also called the nuclear norm or Schatten \(1\)-norm, is
\[
  \norm T_1=\Tr\bigl((T^*T)^{1/2}\bigr),
\]
the sum of the singular values of \(T\).
\end{definition}

Thus for \(G\in\mathcal X=\C^{I\times\Sigma}\),
\[
  \norm{G}_{\HS}^2=\sum_{i\in I}\sum_{\alpha\in\Sigma}|G_{i\alpha}|^2.
\]
The bipartite matrix algebra is \(M_I(\C)\), acting on \(\C^I\).  For
\(A\in M_I(\C)\), the same definitions give
\[
  \norm A_{\op},\qquad \norm A_1,\qquad \norm A_{\HS}.
\]
If \(A\) is Hermitian with eigenvalues \((\lambda_j)\), then these become
\[
  \norm A_{\op}=\max_j|\lambda_j|,\qquad
  \norm A_1=\sum_j|\lambda_j|,\qquad
  \norm A_{\HS}=\left(\sum_j|\lambda_j|^2\right)^{1/2}.
\]

\begin{definition}[Partial transpose]
For \(i=(p,q)\) and \(j=(p',q')\) in \(I=P\times Q\), the partial transpose of
\(A\in M_I(\C)\) along the second tensor factor is
\[
  (A^\pt)_{(p,q),(p',q')}=A_{(p,q'),(p',q)} .
\]
\end{definition}

\begin{lemma}[Elementary norm facts]
For every \(A\in M_I(\C)\),
\[
  \norm{A^\pt}_{\HS}=\norm A_{\HS},\qquad
  \norm A_{\op}\le \norm A_{\HS},\qquad
  \norm A_1\le \sqrt D\,\norm A_{\HS}.
\]
If \(A\) is Hermitian and \(m\) is even, then
\[
  \norm A_{\op}^m\le \Tr(A^m).
\]
\end{lemma}

\begin{proof}
The first identity is only a permutation of matrix coordinates:
\[
  \norm{A^\pt}_{\HS}^2
  =
  \sum_{i,j}|(A^\pt)_{ij}|^2
  =
  \sum_{i,j}|A_{ij}|^2
  =
  \norm A_{\HS}^2 .
\]
Let \(s_1,\ldots,s_D\) be the singular values of \(A\).  By definition,
\[
  \norm A_{\op}=\max_i s_i,\qquad
  \norm A_{\HS}=\left(\sum_i s_i^2\right)^{1/2},\qquad
  \norm A_1=\sum_i s_i .
\]
Therefore
\[
  \norm A_{\op}
  =
  \max_i s_i
  \le
  \left(\sum_i s_i^2\right)^{1/2}
  =
  \norm A_{\HS}.
\]
Cauchy's inequality gives
\[
  \norm A_1
  =
  \sum_i s_i
  \le
  \left(\sum_i1^2\right)^{1/2}
  \left(\sum_i s_i^2\right)^{1/2}
  =
  \sqrt D\,\norm A_{\HS}.
\]
If \(A\) is Hermitian, diagonalize it in an orthonormal basis:
\[
  A=U^*\diag(\lambda_1,\ldots,\lambda_D)U,\qquad \lambda_i\in\R.
\]
Then the singular values are \(|\lambda_i|\), and
\[
  \norm A_{\op}^m=\max_i |\lambda_i|^m\le \sum_i |\lambda_i|^m
  =\sum_i \lambda_i^m=\Tr(A^m),
\]
because \(m\) is even.
\end{proof}

\section{Gaussian, Wishart, density, and direction}

Let \(\gamma\) be the standard complex Gaussian probability measure on
\(\mathcal X\).  Thus all coordinates \(G_{i\alpha}\) are independent standard
complex Gaussians, normalized so that
\[
  \E |G_{i\alpha}|^2=1.
\]
Concretely, a standard complex Gaussian is
\[
  g=\frac{X+iY}{\sqrt2},
\]
where \(X\) and \(Y\) are independent real \(N(0,1)\) variables.  With this
normalization, \(|g|^2\) is an exponential random variable of mean \(1\).
We use the following convention: a Gamma random variable with shape \(r>0\)
and scale \(1\) has density
\[
  \frac{x^{r-1}e^{-x}}{\Gamma(r)}\,\1_{\{x>0\}}.
\]
With this convention, a sum of \(r\) independent exponential variables of mean
\(1\), for integer \(r\), has Gamma distribution with shape \(r\) and scale
\(1\).  Hence
for each fixed \(i\in I\),
\[
  \sum_{\alpha\in\Sigma}|G_{i\alpha}|^2
\]
has the Gamma distribution with shape parameter \(S=\card\Sigma\) and scale
parameter \(1\).  Also
\[
  R(G)^2=\sum_{i\in I}\sum_{\alpha\in\Sigma}|G_{i\alpha}|^2
\]
has Gamma distribution with shape parameter \(DS\) and scale parameter \(1\).
The Gaussian radius and normalized direction are
\[
  R(G)=\norm G_{\HS},\qquad
  \Theta(G)=
  \begin{cases}
    G/\norm G_{\HS},&G\ne 0,\\
    x_0,&G=0,
  \end{cases}
\]
where \(x_0\) is any fixed point of the unit sphere
\[
  \Sphere(\mathcal X)=\set{X\in\mathcal X\given \norm X_{\HS}=1}.
\]
The value at \(0\) is immaterial because the Gaussian has no atom at \(0\).
The Gaussian direction law is
\[
  \nu_{\mathrm{dir}}=\Theta_*\gamma .
\]
The raw Wishart matrix, normalized Wishart matrix, and density matrix are
\[
  \mathcal W_{\mathrm{raw}}(G)=GG^*,\qquad
  W(G)=\frac1S GG^*,\qquad
  \rho(G)=
  \begin{cases}
    GG^*/\Tr(GG^*),&G\ne0,\\
    0,&G=0.
  \end{cases}
\]
Since \(\Tr(GG^*)=R(G)^2\), for \(G\ne0\),
\[
  \rho(G)=\Theta(G)\Theta(G)^* .
\]
The partially transposed density is \(\rho^\pt(G)=(\rho(G))^\pt\).

\begin{lemma}[Radial rescalings]
For \(G\ne0\),
\[
  GG^*=R(G)^2\,\Theta(G)\Theta(G)^*,\qquad
  (GG^*)^\pt=R(G)^2\,(\Theta(G)\Theta(G)^*)^\pt.
\]
Consequently
\[
  \norm{G}_{\op}=R(G)\norm{\Theta(G)}_{\op},\qquad
  \norm{(GG^*)^\pt}_{\op}
  =R(G)^2\norm{(\Theta(G)\Theta(G)^*)^\pt}_{\op}.
\]
\end{lemma}

\begin{proof}
The identities follow from \(G=R(G)\Theta(G)\) and homogeneity of matrix
multiplication, partial transpose, and the operator norm.
\end{proof}

\chapter{Medians, Local Extensions, and Local Levy Reduction}

\section{Medians and center replacement}

Let \((\Omega,\mu)\) be a probability space and \(f:\Omega\to\R\) be
measurable.

\begin{definition}[Median]
A number \(m\in\R\) is a median of \(f\) if
\[
  \mu\{f\le m\}\ge \frac12,\qquad
  \mu\{m\le f\}\ge \frac12.
\]
\end{definition}

\begin{lemma}[Median far from a test center]
Let \(m\) be a median of \(f\).  For every \(a\in\R\) and every \(t>0\), if
\(|m-a|\ge t\), then
\[
  \mu\{|f-a|\ge t\}\ge \frac12.
\]
\end{lemma}

\begin{proof}
If \(m\ge a+t\), then \(\{f\ge m\}\subset\{|f-a|\ge t\}\), and this set has
measure at least \(1/2\).  If \(m\le a-t\), use \(\{f\le m\}\) instead.
\end{proof}

\begin{lemma}[Median tail from any center]
Let \(m\) be a median of \(f\).  For every \(a\in\R\) and \(t>0\),
\[
  \mu\{|f-m|\ge t\}\le 2\,\mu\{|f-a|\ge t/2\}.
\]
\end{lemma}

\begin{proof}
If \(|m-a|<t/2\), then \(|f-m|\ge t\) implies \(|f-a|\ge t/2\).  If
\(|m-a|\ge t/2\), the previous lemma says that the right hand side is at
least \(1\), while the left hand side is a probability and is therefore at
most \(1\).
\end{proof}

\section{Local Lipschitz functions}

\begin{definition}[Local Lipschitz property]
Let \(E\) be a metric space, \(\Omega_0\subset E\), and \(f:E\to\R\).  The
restriction of \(f\) to \(\Omega_0\) is \(L\)-Lipschitz if
\[
  |f(x)-f(y)|\le L\,\dist(x,y)\qquad(x,y\in\Omega_0).
\]
\end{definition}

\begin{lemma}[McShane extension, clipped form]
Assume \(\Omega_0\) is nonempty, \(L,M\ge0\), \(f\) is \(L\)-Lipschitz on
\(\Omega_0\), and \(|f|\le M\) on \(\Omega_0\).  Then there exists an
\(L\)-Lipschitz function
\(\widetilde f:E\to[-M,M]\) such that \(\widetilde f=f\) on \(\Omega_0\).
\end{lemma}

\begin{proof}
Define
\[
  F(x)=\inf_{y\in\Omega_0}\bigl(f(y)+L\,\dist(x,y)\bigr)
\]
for \(x\in E\).  Because \(\Omega_0\) is nonempty and \(|f|\le M\) on
\(\Omega_0\), the infimum is finite.  Indeed, for any fixed \(y_0\in\Omega_0\),
\[
  F(x)\le f(y_0)+L\dist(x,y_0)<\infty,
\]
and for every \(y\in\Omega_0\),
\[
  f(y)+L\dist(x,y)
  \ge
  -M.
\]
Thus \(F(x)\in\R\).

If \(x,x'\in E\), then for each \(y\in\Omega_0\),
\[
  f(y)+L\dist(x,y)
  \le
  f(y)+L\dist(x',y)+L\dist(x,x').
\]
Taking the infimum over \(y\) gives
\[
  F(x)\le F(x')+L\dist(x,x').
\]
Interchanging \(x\) and \(x'\) gives \(|F(x)-F(x')|\le L\dist(x,x')\).

If \(x\in\Omega_0\), then using \(y=x\) gives \(F(x)\le f(x)\).  Conversely,
the \(L\)-Lipschitz condition on \(\Omega_0\) gives
\[
  f(x)\le f(y)+L\dist(x,y)\qquad(y\in\Omega_0),
\]
so \(f(x)\le F(x)\).  Hence \(F=f\) on \(\Omega_0\).

Finally the clipping map \(c(u)=\max(-M,\min(M,u))\) is \(1\)-Lipschitz:
projection onto an interval cannot increase distances.  Therefore
\[
  \widetilde f=c\circ F
\]
is \(L\)-Lipschitz, takes values in \([-M,M]\), and equals \(f\) on
\(\Omega_0\), since \(f(\Omega_0)\subset[-M,M]\).
\end{proof}

\begin{theorem}[Localized Levy reduction]
Let \((S,\dist,\mu)\) be a probability metric space, let \(n>0\), and let
\(L>0\).  Assume that \(S\) satisfies the following global concentration
inequality: for every \(L\)-Lipschitz measurable function \(g:S\to\R\), every
median \(M_g\) of \(g\), and every \(t>0\),
\[
  \mu\{|g-M_g|\ge t\}\le
  2\exp\!\left(-\frac{n t^2}{4L^2}\right)
  .
\]
Let
\(\Omega_0\subset S\) be a nonempty measurable set, let
\(\mu(\Omega_0^c)\le b\), and suppose that \(f:S\to\R\) is measurable and
\(L\)-Lipschitz on \(\Omega_0\).  If \(m_f\) is a median of \(f\), then for
every \(t>0\),
\[
  \mu\{|f-m_f|\ge t\}
  \le
  2b+4\exp\!\left(-\frac{n t^2}{16L^2}\right).
\]
\end{theorem}

\begin{proof}
Extend \(f|_{\Omega_0}\) to an \(L\)-Lipschitz function \(F\) on the whole
space \(S\).  For any median \(M_F\),
\[
  \mu\{|f-M_F|\ge u\}
  \le b+2\exp\!\left(-\frac{n u^2}{4L^2}\right).
\]
Apply the center-replacement lemma with \(u=t/2\):
\[
  \mu\{|f-m_f|\ge t\}
  \le 2\mu\{|f-M_F|\ge t/2\}.
\]
This gives the stated bound.
\end{proof}

\section{Median to mean}

\begin{lemma}[Mean--median gap from an integrated tail]
Let \(f:\Omega\to\R\) be measurable and integrable, let \(m\in\R\), and let
\(R\ge0\).  Assume \(|f-m|\le R\) almost surely.  Let
\(\Phi:[0,R]\to[0,\infty)\) be Borel measurable with finite integral, and
assume
\[
  \mu\{|f-m|\ge t\}\le \Phi(t)\qquad(0\le t\le R).
\]
Then
\[
  |\E f-m|\le\int_0^R \Phi(t)\,\dd t.
\]
\end{lemma}

\begin{proof}
Let \(Y=|f-m|\).  Since \(Y\ge0\),
\[
  Y(\omega)=\int_0^\infty \1_{\{t\le Y(\omega)\}}\,\dd t
\]
for every \(\omega\).  Tonelli's theorem applies because the integrand is
nonnegative, so
\[
  \E|f-m|
  =
  \int_\Omega\int_0^\infty \1_{\{t\le |f-m|\}}\,\dd t\,\dd\mu
  =
  \int_0^\infty \mu\{|f-m|\ge t\}\,\dd t.
\]
The almost sure bound \(|f-m|\le R\) makes the integrand vanish for \(t>R\),
hence
\[
  \E|f-m|
  =
  \int_0^R \mu\{|f-m|\ge t\}\,\dd t
  \le
  \int_0^R\Phi(t)\,\dd t.
\]
Since \(|\E f-m|\le\E|f-m|\), the claim follows.
\end{proof}

\begin{lemma}[Mean deviation from median deviation]
Let \(f:\Omega\to\R\) be measurable and integrable.  If
\(|\E f-m|\le a\), then for every \(u>a\),
\[
  \{|f-\E f|\ge u\}\subset \{|f-m|\ge u-a\}.
\]
In particular, if \(a\le u/2\), then
\[
  \Prob\{|f-\E f|\ge u\}\le \Prob\{|f-m|\ge u/2\}.
\]
\end{lemma}

\begin{proof}
Let \(x\in\Omega\) be such that \(|f(x)-\E f|\ge u\).  By the triangle
inequality,
\[
  |f(x)-m|
  =
  |(f(x)-\E f)+(\E f-m)|
  \ge
  |f(x)-\E f|-|\E f-m|.
\]
Using the hypothesis \(|\E f-m|\le a\), we get
\[
  |f(x)-m|\ge u-a.
\]
This proves
\[
  \{|f-\E f|\ge u\}\subset \{|f-m|\ge u-a\}.
\]
If \(a\le u/2\), then \(u-a\ge u/2\), hence
\[
  \{|f-\E f|\ge u\}\subset \{|f-m|\ge u/2\}.
\]
Taking probabilities gives the final inequality.
\end{proof}

\chapter{Deterministic Matrix Lipschitz Calculus}

\section{Trace powers}

\begin{lemma}[Telescoping identity]
For two elements \(A,B\) of an associative algebra and \(k\ge1\),
\[
  A^k-B^k=\sum_{j=0}^{k-1}A^j(A-B)B^{k-1-j}.
\]
\end{lemma}

\begin{proof}
Expand the right hand side.  Consecutive terms cancel, leaving \(A^k-B^k\).
\end{proof}

\begin{proposition}[Trace-power perturbation]
Let \(A,B\in M_D(\C)\), \(k\ge1\).  Then
\[
  |\Tr(A^k)-\Tr(B^k)|
  \le k\max(\norm A_{\op},\norm B_{\op})^{k-1}\norm{A-B}_1.
\]
\end{proposition}

\begin{proof}
We shall use the finite-dimensional ideal property of the trace norm:
for all composable matrices \(R,S,T\),
\[
  \norm{RST}_1\le \norm R_{\op}\,\norm S_1\,\norm T_{\op}.
\]
Here is the finite-dimensional duality used in the proof.  We claim that
\[
  \norm M_1=\sup_{\norm U_{\op}\le1}|\Tr(UM)|.
\]
For the upper bound, let \(M=V|M|\) be the polar decomposition, where
\(|M|=(M^*M)^{1/2}\).  If \(\norm U_{\op}\le1\), then
\[
  |\Tr(UM)|
  =
  |\Tr(UV|M|)|
  \le
  \norm{UV}_{\op}\Tr|M|
  \le
  \norm M_1.
\]
For the reverse inequality, take \(U=V^*\) on the range of \(|M|\), extended
by \(0\) on its orthogonal complement.  Then \(\norm U_{\op}\le1\) and
\[
  \Tr(UM)=\Tr(V^*V|M|)=\Tr|M|=\norm M_1.
\]
This proves the dual formula.

Therefore
\[
\begin{aligned}
  \norm{RST}_1
  &=
  \sup_{\norm U_{\op}\le1}|\Tr(URST)|  \\
  &=
  \sup_{\norm U_{\op}\le1}|\Tr(TURS)|  \\
  &\le
  \sup_{\norm U_{\op}\le1}\norm{TUR}_{\op}\,\norm S_1 \\
  &\le
  \norm T_{\op}\,\norm R_{\op}\,\norm S_1 .
\end{aligned}
\]
We also use submultiplicativity of the operator norm, so
\(\norm{A^j}_{\op}\le\norm A_{\op}^j\) and
\(\norm{B^{k-1-j}}_{\op}\le\norm B_{\op}^{k-1-j}\).

Now apply the telescoping identity:
\[
  A^k-B^k=\sum_{j=0}^{k-1}M_j,
  \qquad
  M_j=A^j(A-B)B^{k-1-j}.
\]
By linearity of the trace, the triangle inequality, and
\(|\Tr M|\le\norm M_1\), we have
\[
\begin{aligned}
|\Tr(A^k)-\Tr(B^k)|
&=
\left|\sum_{j=0}^{k-1}\Tr(M_j)\right|\\
&\le
\sum_{j=0}^{k-1}|\Tr(M_j)|\\
&\le
\sum_{j=0}^{k-1}\norm{M_j}_1.
\end{aligned}
\]
Therefore
\[
\begin{aligned}
|\Tr(A^k)-\Tr(B^k)|
&\le \sum_{j=0}^{k-1}
  \norm{A^j(A-B)B^{k-1-j}}_1\\
&\le \sum_{j=0}^{k-1}
  \norm A_{\op}^j\norm{A-B}_1\norm B_{\op}^{k-1-j}\\
&\le k\max(\norm A_{\op},\norm B_{\op})^{k-1}\norm{A-B}_1.
\end{aligned}
\]
\end{proof}

\begin{proposition}[Partial-transpose moment Lipschitz bound]
Let \(k\ge1\).  For \(X,Y\in\Sphere(\mathcal X)\), put
\[
  A_X=(XX^*)^\pt,\qquad A_Y=(YY^*)^\pt .
\]
Then
\[
\begin{aligned}
&|\Tr(A_X^k)-\Tr(A_Y^k)|\\
&\quad\le
  k\,D^{1/2}\,
  \max(\norm{A_X}_{\op},\norm{A_Y}_{\op})^{k-1}
  (\norm X_{\op}+\norm Y_{\op})\,\norm{X-Y}_{\HS}.
\end{aligned}
\]
\end{proposition}

\begin{proof}
The trace-power perturbation gives
\[
  |\Tr(A_X^k)-\Tr(A_Y^k)|
  \le
  k\max(\norm{A_X}_{\op},\norm{A_Y}_{\op})^{k-1}\norm{A_X-A_Y}_1.
\]
Since partial transpose is Hilbert--Schmidt isometric,
\[
  \norm{A_X-A_Y}_1
  \le D^{1/2}\norm{A_X-A_Y}_{\HS}
  =D^{1/2}\norm{XX^*-YY^*}_{\HS}.
\]
Now
\[
  XX^*-YY^*=(X-Y)X^*+Y(X-Y)^*,
\]
and the Hilbert--Schmidt ideal inequality gives
\[
  \norm{XX^*-YY^*}_{\HS}
  \le(\norm X_{\op}+\norm Y_{\op})\norm{X-Y}_{\HS}.
\]
\end{proof}

\begin{corollary}[Good-set Lipschitz scale]
Let \(k\ge1\), \(a,b\ge0\), and \(d\ge1\).  Suppose that for every
\(X\in\Omega_0\subset\Sphere(\mathcal X)\) one has
\[
  \norm X_{\op}\le \frac a d,\qquad
  \norm{(XX^*)^\pt}_{\op}\le \frac b{d^2},
\]
and \(D=d^2\), then
\[
  X\mapsto \Tr\left(\left((XX^*)^\pt\right)^k\right)
\]
is \(L\)-Lipschitz on \(\Omega_0\), with
\[
  L\le 2kab^{k-1}d^{-(2k-2)}.
\]
\end{corollary}

\begin{proof}
Substitute the two good-set bounds and \(D^{1/2}=d\) into the preceding
proposition.
\end{proof}

\chapter{Complex Gaussian Wick Expansion}

\section{Monomials and pairings}

Let \(\Xi\) be a finite set, and let \(\{z_\eta\}_{\eta\in\Xi}\) be
independent standard complex Gaussian variables.  A monomial is encoded by
two finite words
\[
  h=(h_0,\ldots,h_{r-1}),\qquad
  \bar h=(\bar h_0,\ldots,\bar h_{s-1})
\]
in \(\Xi\), representing
\[
  z_{h_0}\cdots z_{h_{r-1}}
  \overline{z_{\bar h_0}}\cdots\overline{z_{\bar h_{s-1}}}.
\]

\begin{theorem}[Complex Wick formula]
The expectation of the monomial is zero unless \(r=s\).  If \(r=s\), then
\[
  \E\prod_{a=0}^{r-1}z_{h_a}
    \prod_{b=0}^{r-1}\overline{z_{\bar h_b}}
  =
  \card\{\pi\in\mathfrak S_r:\ h_a=\bar h_{\pi(a)}
       \text{ for every }a\}.
\]
\end{theorem}

\begin{proof}
Group the factors by coordinate label.  Since the coordinates are
independent, the expectation factors over labels \(\eta\in J\).  Suppose that
the label \(\eta\) appears \(a_\eta\) times among
\(h_0,\ldots,h_{r-1}\) and \(b_\eta\) times among
\(\bar h_0,\ldots,\bar h_{r-1}\).  The \(\eta\)-factor is
\[
  \E z_\eta^{a_\eta}\overline{z_\eta}^{\,b_\eta}.
\]
Write \(z_\eta=Re^{i\Theta}\) in polar coordinates.  The angle \(\Theta\) is
uniform on \([0,2\pi]\), so this expectation contains the angular integral
\[
  \frac1{2\pi}\int_0^{2\pi} e^{i(a_\eta-b_\eta)\theta}\,\dd\theta.
\]
This integral is \(0\) unless \(a_\eta=b_\eta\).  If
\(a_\eta=b_\eta=s_\eta\), the same coordinate contributes
\[
  \E |z_\eta|^{2s_\eta}.
\]
For a standard complex Gaussian, \(|z_\eta|^2\) has density \(e^{-x}\1_{x\ge0}\),
so
\[
  \E |z_\eta|^{2s_\eta}
  =
  \int_0^\infty x^{s_\eta}e^{-x}\,\dd x
  =
  s_\eta!.
\]
Thus the whole expectation is zero unless every coordinate appears equally
often holomorphically and anti-holomorphically.  In the nonzero case it is
\[
  \prod_{\eta\in J}s_\eta!.
\]
But \(\prod_\eta s_\eta!\) is exactly the number of bijections from the
holomorphic positions to the anti-holomorphic positions that preserve the
coordinate label: for each label \(\eta\), one may match the
\(s_\eta\) holomorphic occurrences with the \(s_\eta\) anti-holomorphic
occurrences in \(s_\eta!\) ways, independently over \(\eta\).  This is the
cardinality displayed in the theorem.
\end{proof}

\section{Closed walks and surviving contractions}

For \(m\ge0\), a closed walk of length \(m+1\) in \(I\) is a tuple
\[
  w=(i_0,i_1,\ldots,i_m,i_{m+1}),\qquad i_{m+1}=i_0.
\]
Equivalently, it is a base point and a word in \(I\).  A sample word is
\(\alpha:\{0,\ldots,m\}\to\Sigma\).

Each edge \(e\) contributes a holomorphic Gaussian coordinate and an
anti-holomorphic Gaussian coordinate.  Because the matrix has been partially
transposed, the two coordinates use crossed \(P,Q\)-components.
Writing \(i_e=(p_e,q_e)\), these two coordinates are
\[
  (p_e,q_{e+1},\alpha(e)),
  \qquad
  (p_{e+1},q_e,\alpha(e)).
\]

\begin{definition}[Wick permutation]
For a monomial containing \(m+1\) holomorphic Gaussian factors and \(m+1\)
anti-holomorphic Gaussian factors, a Wick permutation is a permutation
\(\pi\in\mathfrak S_{m+1}\).  It records the Wick contraction in which the
holomorphic factor at edge \(e\) is paired with the anti-holomorphic factor at
edge \(\pi(e)\).
\end{definition}

A Wick permutation survives a given closed walk and sample word if every
holomorphic edge coordinate at \(e\) equals the anti-holomorphic edge
coordinate at \(\pi(e)\).

\begin{definition}[Surviving closed-walk pairing]
For a closed walk \(w\) and sample word \(\alpha\), the set
\(\mathcal S(w,\alpha)\) consists of all permutations \(\pi\in\mathfrak S_{m+1}\)
such that, for every edge \(e\),
\[
  (p_e,q_{e+1},\alpha(e))
  =
  (p_{\pi(e)+1},q_{\pi(e)},\alpha(\pi(e))),
\]
with edge indices read cyclically.  Equivalently, the holomorphic coordinate
on edge \(e\) is exactly the anti-holomorphic coordinate on edge \(\pi(e)\).
\end{definition}

\begin{proposition}[Moment as a surviving-pairing sum]
Let \(Z=W^\pt-\diag(W^\pt)\).  For every \(m\),
\[
  \E\Tr(Z^{m+1})
  =
  \sum_{w}\sum_{\alpha}
  c(w,\alpha)\,\card\mathcal S(w,\alpha),
\]
where \(w\) ranges over closed walks of length \(m+1\), \(\alpha\) ranges over
sample words, and \(c(w,\alpha)\) is the product of the off-diagonal edge
coefficients.
\end{proposition}

\begin{proof}
Expand the trace as a sum over closed walks, expand each matrix entry in
sample coordinates, then apply the complex Wick formula to each monomial.  The
off-diagonal projection kills any path containing a diagonal edge, and the
surviving Wick permutations are exactly the permutations counted by
\(\mathcal S(w,\alpha)\).
\end{proof}

\chapter{Aubrun's Combinatorial Counting Layer}

\section{Equivalence relations associated with a Wick permutation}

Fix \(m\) and \(\pi\in\mathfrak S_{m+1}\).  The survival constraints generate
three equivalence relations:
\[
  \sim_L,\qquad \sim_R,\qquad \sim_\Sigma
\]
on the edge positions.  They identify, respectively, left coordinates, right
coordinates, and sample-column coordinates that must agree for \(\pi\) to
survive.  Let
\[
  \ell(\pi),\quad r(\pi),\quad c(\pi)
\]
be the numbers of equivalence classes.  The weighted rank is
\[
  \rho(\pi)=\ell(\pi)+r(\pi)+c(\pi),
\]
and the defect is the gap between the naive \(3(m+1)\) degrees of freedom and
the surviving degrees of freedom after the Wick equalities.

For this section, the fiber \(\mathcal F_\pi\) is the set of pairs
\((w,\alpha)\) such that the Wick permutation \(\pi\) survives the closed walk
\(w\) and sample word \(\alpha\).

\begin{proposition}[Fiber cardinality bound]
For a fixed Wick permutation \(\pi\), the number of triples
\((w,\alpha,\pi)\) satisfying the survival constraints is bounded by the
product of the powers of the three ambient cardinalities attached to the class
counts:
\[
  \card \mathcal F_\pi
  \le
  |P|^{\ell(\pi)}|Q|^{r(\pi)}|\Sigma|^{c(\pi)}.
\]
\end{proposition}

\begin{proof}
Choose one coordinate for each equivalence class.  All other coordinates are
then forced by the relation.  The left, right, and sample coordinates may be
chosen independently from \(P,Q,\Sigma\), respectively.
\end{proof}

\section{Tilde words and Aubrun's lists}

The original paper reorganizes the same data into three words
\[
  \widetilde a,\qquad \widetilde b,\qquad \widetilde c
\]
and two lists, conventionally called list (3) and list (7).  In the formal
development these are abstracted as follows.

\begin{definition}[Aubrun tilde words]
Let \(A,B,C\) be finite sets.  An Aubrun word of length \(m+1\) is a triple
\[
  (\widetilde a,\widetilde b,\widetilde c)
  \in A^{m+1}\times B^{m+1}\times C^{m+1}.
\]
List (3) is the edge list
\[
  ((\widetilde a_e,\widetilde b_e,\widetilde c_e))_{e=0}^m,
\]
and list (7) is the Wick-matched list obtained by applying the Wick
permutation and the partial-transpose crossing to the same positions.
\end{definition}

\begin{definition}[Matching and nonrepetition conditions]
The Wishart matching condition says that the column labels on matched
holomorphic and anti-holomorphic variables agree.  The triple matching
condition says that the left, right, and column labels agree in the crossed
pattern dictated by partial transpose.  The pair
\((\widetilde a,\widetilde b)\) is nonrepetitive if no adjacent edge repeats
the same ordered pair; this records the off-diagonal projection.
\end{definition}

\begin{definition}[Innovations and compatibility]
An innovation is a first occurrence of a new equivalence class while scanning
the word.  The compatibility defect counts how many apparent choices are lost
when the \(A\)- and \(C\)-relations are required to be simultaneously
compatible with the Wick relation.  If \(U\) denotes the number of innovations,
then the relevant class-count loss is at least the compatibility defect.
\end{definition}

A defect profile records the finite data used in this last count: the number
\(U\) of innovations, the defect \(\Delta\), and the resulting class counts.
For fixed \(m\), there are only finitely many possible defect profiles.

\begin{lemma}[Aubrun Lemma 3.6]
For fixed innovation pattern with \(U\) innovations, the number of compatible
tilde-word assignments is at most
\[
  k^{3U}.
\]
\end{lemma}

\begin{proof}
At each innovation one may choose at most one new value in each of the three
coordinate alphabets \(A,B,C\).  Each such choice has at most \(k\) options in
the canonical finite model.  Noninnovative positions are forced by previous
equivalence-class representatives and the matching constraints.  Thus every
innovation contributes at most \(k^3\), giving \(k^{3U}\).
\end{proof}

\begin{lemma}[Aubrun Lemma 7.4]
Fix the word length \(m+1\), an innovation count \(U\), and a defect
\(\Delta\).  The number \(N_\Delta\) of compatible equivalence-class data with
that defect is bounded by
\[
  N_{\Delta}\le C(m)\,k^{3U-\Delta}.
\]
Here \(k\) is the common upper bound for the finite alphabets being chosen
from, and \(C(m)\) depends only on the word length.
\end{lemma}

\begin{proof}
Lift the innovation set from \(m\) edge positions to \(m+1\) vertex positions.
There are only \(C_1(m)\) ways to choose the finite combinatorial pattern
recording where the innovations and forced identifications occur.  Once that
pattern is fixed, Lemma 3.6 gives at most \(k^{3U}\) raw choices.  A defect
\(\Delta\) means that \(\Delta\) of these apparent choices are not genuinely
free: each such lost degree of freedom replaces one free \(k\)-choice by a
forced equality with a previously chosen representative.  Therefore the
choice count for the fixed pattern is at most \(k^{3U-\Delta}\).  Multiplying
by the number \(C_1(m)\) of possible finite patterns gives the asserted
bound, with \(C(m)=C_1(m)\).
\end{proof}

\begin{lemma}[Aubrun Lemma 7.5]
For each fixed \(m\), there is a polynomial \(P_m\), depending only on \(m\),
such that the total surviving relation count satisfies
\[
  \sum_{\pi\in\mathfrak S_{m+1}}\card\mathcal F_\pi
  \le P_m(k).
\]
\end{lemma}

\begin{proof}
Partition the Wick permutations by their defect profile.  Lemma 7.4 bounds
each profile class by \(C(m)k^{3U-\Delta}\).  For fixed \(m\), the parameters
\(U\) and \(\Delta\), and the auxiliary equivalence-relation pattern, range
over a finite set depending only on \(m\).  Hence the sum of all profile
bounds is a finite sum of monomials in \(k\), with coefficients depending only
on \(m\).  This finite sum is a polynomial \(P_m(k)\).  The profile-sum
identity identifies the sum of the profile counts with
\(\sum_{\pi\in\mathfrak S_{m+1}}\card\mathcal F_\pi\), proving the displayed
bound.
\end{proof}

\begin{theorem}[Aubrun relation-count bound]
Fix \(m\ge0\).  For \(u,v\in\N\), let \(N_m(u,v)\) be the number of surviving
Wick permutations \(\pi\in\mathfrak S_{m+1}\) for which
\[
  u=\ell(\pi)+r(\pi),
  \qquad
  v=c(\pi).
\]
Thus \(u\) counts the free left/right classes and \(v\) counts the free
sample-column classes.  There is a polynomial \(Q\), depending only on the
Wick length and not on the dimensions \(d\) and \(s\), such that for all
\(d,s\ge1\),
\[
  \sum_{u,v} N_m(u,v)\,d^u\,s^v
  \le
  2^{m+1}
  \bigl(d+Q(m+1)\bigr)^{m+3}
  \bigl(\sqrt s+Q(m+1)\bigr)^{m+1}.
\]
\end{theorem}

\begin{proof}
After Wick's formula, the moment is a finite sum over surviving permutations.
For a permutation with profile \((u,v)\), the fiber-cardinality bound above
shows that its contribution is controlled by \(d^u s^v\): one chooses one
left/right coordinate for each of the \(u\) free left/right classes and one
sample label for each of the \(v\) free column classes.

It remains to count how many permutations have each profile.  The tilde-word
analysis does this by grouping permutations according to their defect.  Lemma
3.6 bounds the choices made at innovation positions.  Lemma 7.4 bounds the
number of compatible equivalence classes at fixed defect.  Lemma 7.5 then sums
over the finitely many defect profiles.  The resulting correction is absorbed
into the polynomial \(Q\), which depends only on \(m\).  Substituting this
profile count into the weighted profile sum gives the displayed bound.
\end{proof}

\chapter{Aubrun Moment Inputs}

\section{Off-diagonal moment envelope}

Let
\[
  Z=W^\pt-\diag(W^\pt).
\]
The off-diagonal trace moment is
\[
  M_m=\E\Tr(Z^m).
\]
The Wick expansion and the relation-count bound from the previous chapter
imply a concrete envelope
\[
  |M_m|
  \le
  \left(\frac2S\right)^m
  (d+Q(m))^{m+2}(\sqrt S+Q(m))^m,
\]
for a polynomial function \(Q\) independent of \(d\) and \(S\).

\begin{theorem}[Aubrun off-diagonal moment bound]
For every \(m\ge1\), there is a polynomial \(Q\), independent of \(d\) and
\(S\), such that the off-diagonal partially transposed Wishart moment satisfies
\[
  |M_m|
  \le
  \left(\frac2S\right)^m
  (d+Q(m))^{m+2}(\sqrt S+Q(m))^m.
\]
\end{theorem}

\begin{proof}
Write a cyclic word in \(I=P\times Q\) as
\[
  w=(i_0,i_1,\ldots,i_m),\qquad i_m=i_0,
\]
and write \(i_r=(p_r,q_r)\).  Expanding the trace gives
\[
  \Tr(Z^m)
  =
  \sum_w\prod_{r=0}^{m-1}Z_{i_r,i_{r+1}}.
\]
For an off-diagonal edge \(i_r\ne i_{r+1}\), the partial-transposed Wishart
entry is
\[
  Z_{i_r,i_{r+1}}
  =
  \frac1S
  \sum_{\alpha_r\in\Sigma}
  G_{(p_r,q_{r+1}),\alpha_r}\,
  \overline{G_{(p_{r+1},q_r),\alpha_r}}.
\]
For a diagonal edge \(i_r=i_{r+1}\), the entry of \(Z\) is zero.  Therefore
\[
  M_m
  =
  S^{-m}
  \sum_w\mathbf 1_{\mathrm{off}}(w)
  \sum_{\alpha_0,\ldots,\alpha_{m-1}\in\Sigma}
  \E\prod_{r=0}^{m-1}
  G_{(p_r,q_{r+1}),\alpha_r}\,
  \overline{G_{(p_{r+1},q_r),\alpha_r}}.
\]
By the complex Wick formula, the expectation is the number of Wick
permutations \(\pi\in\mathfrak S_m\) such that, for every \(r\),
\[
  (p_r,q_{r+1},\alpha_r)
  =
  (p_{\pi(r)+1},q_{\pi(r)},\alpha_{\pi(r)}),
\]
with indices read cyclically.  Thus
\[
  M_m
  =
  S^{-m}
  \sum_{\pi\in\mathfrak S_m}
  \sum_{(w,\alpha)\in\mathcal F_\pi}
  \mathbf 1_{\mathrm{off}}(w).
\]
Since \(0\le \mathbf 1_{\mathrm{off}}\le1\),
\[
  |M_m|
  \le
  S^{-m}\sum_{\pi\in\mathfrak S_m}\card\mathcal F_\pi.
\]
If \(\pi\) has profile \((u,v)\), the fiber-cardinality bound gives
\(\card\mathcal F_\pi\le d^uS^v\).  Regrouping by profile,
\[
  \sum_{\pi\in\mathfrak S_m}\card\mathcal F_\pi
  \le
  \sum_{u,v}N_m(u,v)d^uS^v.
\]
The relation-count bound from the previous chapter gives
\[
  \sum_{u,v}N_m(u,v)d^uS^v
  \le
  2^m(d+Q(m))^{m+2}(\sqrt S+Q(m))^m.
\]
Multiplying by \(S^{-m}\) yields
\[
  |M_m|
  \le
  \left(\frac2S\right)^m
  (d+Q(m))^{m+2}(\sqrt S+Q(m))^m.
\]
\end{proof}

\section{From moments to operator norm expectation}

\begin{lemma}[Even moment to operator norm]
Let \(Z\) be a random Hermitian matrix and let \(m\ge2\) be even.  Assume
\(\E\Tr(Z^m)<\infty\).  Then
\[
  \E\norm Z_{\op}\le \bigl(\E\Tr(Z^m)\bigr)^{1/m}.
\]
\end{lemma}

\begin{proof}
Pointwise, \(\norm Z_{\op}^m\le\Tr(Z^m)\).  Therefore
\[
  \E\norm Z_{\op}^m\le \E\Tr(Z^m).
\]
Since the underlying measure is a probability measure, Holder's inequality
with exponents \(m\) and \(m/(m-1)\) gives
\[
  \E\norm Z_{\op}
  =
  \E(\norm Z_{\op}\cdot 1)
  \le
  \bigl(\E\norm Z_{\op}^m\bigr)^{1/m}(\E 1)^{(m-1)/m}
  =
  \bigl(\E\norm Z_{\op}^m\bigr)^{1/m},
\]
and the two inequalities combine.
\end{proof}

\begin{theorem}[Off-diagonal expectation bound]
Fix \(\lambda>0\) and assume \(D=d^2\) and \(S=s_d\) with
\(s_d/d^2\to\lambda\).  Then there exists \(C_\lambda\) such that
\[
  \E\norm{W^\pt-\diag(W^\pt)}_{\op}\le C_\lambda .
\]
\end{theorem}

\begin{proof}
Choose an even \(m=m(d)\) tending to infinity slowly enough that
\[
  Q(m)=o(d),\qquad \frac{\log d}{m}\to0.
\]
Apply the even-moment lemma to \(Z=W^\pt-\diag(W^\pt)\), and then insert the
moment envelope:
\[
  \E\norm Z_{\op}
  \le \bigl(\E\Tr(Z^m)\bigr)^{1/m}
  \le
  \frac2S(d+Q(m))^{1+2/m}(\sqrt S+Q(m)).
\]
Rewrite the right hand side as
\[
  2\,\frac d{\sqrt S}\,
  d^{2/m}
  \left(1+\frac{Q(m)}d\right)^{1+2/m}
  \left(1+\frac{Q(m)}{\sqrt S}\right).
\]
Because \(S\sim\lambda d^2\), the first factor \(d/\sqrt S\) tends to
\(\lambda^{-1/2}\).  The condition \(\log d/m\to0\) gives
\(d^{2/m}=\exp(2\log d/m)\to1\).  The condition \(Q(m)=o(d)\) gives
\(Q(m)/d\to0\), and also \(Q(m)/\sqrt S\to0\).  Hence the displayed product is
bounded uniformly in \(d\).
\end{proof}

\chapter{Diagonal Gamma and Wishart Expectations}

\section{Diagonal part}

The diagonal entries of \(W^\pt\) are the diagonal entries of \(W\):
\[
  (W^\pt)_{ii}=W_{ii}=\frac1S\sum_{\alpha\in\Sigma}|G_{i\alpha}|^2.
\]
Thus the diagonal operator norm is the maximum of \(D\) averages of \(S\)
independent exponential variables.

\begin{lemma}[Gamma maximum tail]
Assume \(D\ge1\) and \(S\ge1\).  For every \(u\ge0\),
\[
  \Prob\left\{
    \max_{i\in I}W_{ii}
    \ge 1+2\sqrt{\frac{u+\log D}{S}}
       +2\frac{u+\log D}{S}
  \right\}\le e^{-u}.
\]
\end{lemma}

\begin{proof}
First prove the one-coordinate tail.  Let
\[
  \xi=\frac1S\sum_{\alpha=1}^S E_\alpha,
\]
where the \(E_\alpha\) are independent exponential variables of mean \(1\).
For \(0\le\theta<1\),
\[
  \E e^{\theta S\xi}
  =
  \prod_{\alpha=1}^S\E e^{\theta E_\alpha}
  =
  (1-\theta)^{-S},
\]
because \(\E e^{\theta E_\alpha}=\int_0^\infty e^{\theta x}e^{-x}\,\dd x
=(1-\theta)^{-1}\).  Put \(y=\sqrt{v/S}\).  If \(v=0\), the desired
one-coordinate estimate is just a probability bound by \(1\).  If \(v>0\),
choose \(\theta=y/(1+y)\).  Markov's inequality gives
\[
  \Prob\{\xi\ge1+2y+2y^2\}
  \le
  \exp\{-S\theta(1+2y+2y^2)\}(1-\theta)^{-S}.
\]
Since \(1-\theta=(1+y)^{-1}\), the exponent equals
\[
  -S\bigl(\theta(1+2y+2y^2)-\log(1+y)\bigr).
\]
Now
\[
  \theta(1+2y+2y^2)-y^2
  =
  y+\frac{y^3}{1+y}
  \ge y
  \ge\log(1+y),
\]
so the exponent is at most \(-Sy^2=-v\).  Hence
\[
  \Prob\left\{\xi\ge1+2\sqrt{v/S}+2v/S\right\}\le e^{-v}.
\]
Apply this with \(v=u+\log D\).  For each fixed \(i\),
\[
  \Prob\left\{W_{ii}\ge
    1+2\sqrt{\frac{u+\log D}{S}}
       +2\frac{u+\log D}{S}
  \right\}
  \le e^{-u-\log D}=\frac{e^{-u}}{D}.
\]
The event that the maximum exceeds the displayed threshold is contained in
the union of these \(D\) one-coordinate events.  The finite union bound gives
\[
  \Prob\left\{\max_{i\in I}W_{ii}\ge
    1+2\sqrt{\frac{u+\log D}{S}}
       +2\frac{u+\log D}{S}
  \right\}
  \le
  D\cdot\frac{e^{-u}}{D}=e^{-u}.
\]
\end{proof}

\begin{theorem}[Diagonal expectation bound]
Fix \(\lambda>0\) and assume \(D=d^2\) and \(S=s_d\) with
\(s_d/d^2\to\lambda\).  Then
\[
  \E\norm{\diag(W^\pt)}_{\op}\le C_\lambda.
\]
\end{theorem}

\begin{proof}
Let
\[
  M=\norm{\diag(W^\pt)}_{\op}=\max_{i\in I}W_{ii}
\]
and set
\[
  a=1+2\sqrt{\frac{\log D}{S}}+2\frac{\log D}{S}.
\]
Since \(\sqrt{u+\log D}\le \sqrt u+\sqrt{\log D}\), the previous lemma gives,
for every \(u\ge0\),
\[
  \Prob\left\{
    M\ge a+2\sqrt{\frac uS}+2\frac uS
  \right\}
  \le e^{-u}.
\]
Put
\[
  h(u)=2\sqrt{\frac uS}+2\frac uS.
\]
For the nonnegative random variable \((M-a)_+\), Tonelli's theorem gives
\[
  (M-a)_+=\int_0^\infty \1_{\{y\le (M-a)_+\}}\,\dd y
\]
pointwise, and therefore
\[
  \E(M-a)_+=\int_0^\infty \Prob\{(M-a)_+\ge y\}\,\dd y .
\]
Changing variables \(y=h(u)\), which is increasing on \([0,\infty)\), yields
\[
  \E(M-a)_+
  \le
  \int_0^\infty e^{-u}h'(u)\,\dd u
  =
  \frac1{\sqrt S}\int_0^\infty u^{-1/2}e^{-u}\,\dd u
  +\frac2S\int_0^\infty e^{-u}\,\dd u .
\]
The two integrals are \(\Gamma(1/2)=\sqrt\pi\) and \(1\), respectively, so
\[
  \E M
  \le
  1+2\sqrt{\frac{\log D}{S}}+2\frac{\log D}{S}
  +\frac{\sqrt\pi}{\sqrt S}+\frac2S.
\]
In the balanced regime \(D=d^2\) and \(S\sim\lambda d^2\), the right hand side
is bounded uniformly in \(d\); indeed all correction terms tend to \(0\).
\end{proof}

\section{Full expectation}

\begin{theorem}[Wishart-gamma expectation bound]
Fix \(\lambda>0\) and assume \(D=d^2\) and \(S=s_d\) with
\(s_d/d^2\to\lambda\).  Then
\[
  \E\norm{W^\pt}_{\op}\le C_\lambda.
\]
\end{theorem}

\begin{proof}
Use the decomposition
\[
  W^\pt=\diag(W^\pt)+(W^\pt-\diag(W^\pt))
\]
and the operator-norm triangle inequality:
\[
  \norm{W^\pt}_{\op}
  \le
  \norm{\diag(W^\pt)}_{\op}
  +
  \norm{W^\pt-\diag(W^\pt)}_{\op}.
\]
Taking expectations and applying the diagonal and off-diagonal expectation
bounds gives
\[
  \E\norm{W^\pt}_{\op}
  \le C_\lambda^{\mathrm{diag}}+C_\lambda^{\mathrm{off}}.
\]
Rename the sum of the two constants \(C_\lambda\).
\end{proof}

\chapter{High-Probability Bounds}

\section{Gaussian quadratic Bernstein}

Let \(g\) be a standard complex Gaussian vector in \(\C^n\), and let \(A\) be
a fixed Hermitian matrix.  The centered quadratic form
\[
  \ip{Ag}{g}-\E\ip{Ag}{g}
\]
has sub-exponential tails governed by \(\norm A_{\HS}\) and
\(\norm A_{\op}\).

\begin{theorem}[Quadratic-form Bernstein]
Let \(A=A^*\in M_n(\C)\).  If \(A\ne0\), then, for every \(t\ge0\),
\[
  \Prob\{|\ip{Ag}{g}-\Tr A|\ge t\}
  \le
  2\exp\left(
    -\frac18\min\left(\frac{t^2}{\norm A_{\HS}^2},
                \frac{t}{\norm A_{\op}}\right)
  \right).
\]
If \(A=0\), the centered quadratic form is identically zero.
\end{theorem}

\begin{proof}
We prove the estimate directly from the scalar moment-generating functions.

The case \(t=0\) is immediate, so assume \(t>0\).  Diagonalize the Hermitian
matrix \(A\) in an orthonormal basis: choose a unitary \(U\) and real
eigenvalues \(\lambda_1,\ldots,\lambda_n\) such that
\[
  A=U^*\diag(\lambda_1,\ldots,\lambda_n)U .
\]
The standard complex Gaussian law has density \(\pi^{-n}e^{-\norm z_2^2}\)
with respect to Lebesgue measure on \(\C^n\simeq\R^{2n}\).  Since \(U\)
preserves \(\norm z_2\) and has real Jacobian \(1\), the vector
\(\zeta=Ug\) is again standard complex Gaussian.  Its coordinates are
independent standard complex Gaussians, and therefore
\[
  X:=\ip{Ag}{g}-\Tr A
   =\sum_{j=1}^n \lambda_j\bigl(|\zeta_j|^2-1\bigr).
\]
For one standard complex Gaussian coordinate \(\zeta_j\), the variable
\(R_j=|\zeta_j|^2\) is exponential with mean \(1\).  Hence, for every
real \(b<1\),
\[
  \E e^{b(R_j-1)}
  = e^{-b}\int_0^\infty e^{br}e^{-r}\,\dd r
  = \frac{e^{-b}}{1-b}.
\]
Consequently, if \(\theta\in\R\) satisfies
\(|\theta|\le(2\norm A_{\op})^{-1}\), then
\(|\theta\lambda_j|\le1/2\) for all \(j\), and independence gives
\[
  \E e^{\theta X}
  =
  \prod_{j=1}^n
  \frac{\exp(-\theta\lambda_j)}{1-\theta\lambda_j}.
\]
We now bound each scalar factor.  The elementary inequality
\(\log y\le y-1\) for \(y>0\) follows by applying calculus to
\(y-1-\log y\), whose minimum is \(0\) at \(y=1\).  With
\(y=(1-b)^{-1}\) and \(|b|\le1/2\), this gives
\[
  -\log(1-b)\le \frac{b}{1-b},
\]
and hence
\[
  -b-\log(1-b)
  \le -b+\frac{b}{1-b}
  =\frac{b^2}{1-b}
  \le 2b^2.
\]
Applying this with \(b=\theta\lambda_j\) and multiplying the factors,
\[
  \E e^{\theta X}
  \le
  \prod_{j=1}^n \exp\bigl(2\theta^2\lambda_j^2\bigr)
  =
  \exp\bigl(2\theta^2\norm A_{\HS}^2\bigr).
\]

Write
\[
  V=\norm A_{\HS}^2,\qquad L=\norm A_{\op}.
\]
Since \(A\ne0\), both \(V\) and \(L\) are positive.  Markov's inequality,
applied to the nonnegative random variable \(e^{\theta X}\), says that for
\(\theta\ge0\),
\[
  \Prob\{X\ge t\}
  =
  \Prob\{e^{\theta X}\ge e^{\theta t}\}
  \le e^{-\theta t}\E e^{\theta X}.
\]
Choose
\[
  \theta=\min\left(\frac{t}{4V},\frac1{2L}\right).
\]
Then \(\theta\ge0\) and \(\theta\le(2L)^{-1}\), so the MGF bound applies and
\[
  \Prob\{X\ge t\}
  \le
  \exp\{-\theta t+2\theta^2V\}.
\]
It remains only to record the elementary optimization.  If
\(\theta=t/(4V)\), then
\[
  \theta t-2\theta^2V=\frac{t^2}{8V}
  \ge \frac18\min\left(\frac{t^2}{V},\frac{t}{L}\right).
\]
If \(\theta=(2L)^{-1}\), then \(t/(4V)\ge(2L)^{-1}\), hence
\(tL\ge2V\), and
\[
  \theta t-2\theta^2V
  =\frac{t}{2L}-\frac{V}{2L^2}
  \ge \frac{t}{4L}
  \ge \frac18\min\left(\frac{t^2}{V},\frac{t}{L}\right).
\]
Thus
\[
  \Prob\{X\ge t\}
  \le
  \exp\left\{
    -\frac18\min\left(\frac{t^2}{\norm A_{\HS}^2},
                      \frac{t}{\norm A_{\op}}\right)
  \right\}.
\]
The same argument applied to \(-X\) gives the identical bound for
\(\Prob\{X\le -t\}\), because the MGF estimate was proved for positive and
negative \(\theta\) in the same range.  Finally,
\[
  \{|X|\ge t\}\subset \{X\ge t\}\cup\{X\le -t\},
\]
and the union bound gives the displayed factor \(2\).
\end{proof}

\section{Nets and operator norms}

\begin{lemma}[Quadratic net lifting]
Let \(T\) be a Hermitian operator on a Hilbert space, and let \(N\) be a
\(1/4\)-net of the unit sphere.  If
\[
  |\ip{Tu}{u}|\le B\qquad(u\in N),
\]
then
\[
  \norm T_{\op}\le 2B.
\]
\end{lemma}

\begin{proof}
For any unit vector \(x\), choose \(u\in N\) with \(\norm{x-u}\le1/4\).  The
identity
\[
  \ip{Tx}{x}-\ip{Tu}{u}
  =
  \ip{T(x-u)}{x}+\ip{Tu}{x-u}
\]
and Cauchy's inequality give
\[
  |\ip{Tx}{x}-\ip{Tu}{u}|
  \le
  \norm T_{\op}\norm{x-u}\norm x
  +
  \norm T_{\op}\norm u\norm{x-u}
  =
  2\norm T_{\op}\norm{x-u}.
\]
Since \(T\) is Hermitian,
\[
  \norm T_{\op}=\sup_{\norm x=1}|\ip{Tx}{x}|.
\]
Taking the supremum over unit \(x\) gives
\[
  \norm T_{\op}\le B+\frac12\norm T_{\op}.
\]
Moving the last term to the left gives \(\norm T_{\op}\le2B\).
\end{proof}

\begin{lemma}[Net cardinality]
The unit sphere in a real vector space of dimension \(n\) has a \(1/4\)-net
with at most \(9^n\) points.  In the complex \(N\)-dimensional case this is
at most \(81^N\), and \(81^N\le e^{7N}\).
\end{lemma}

\begin{proof}
Take a maximal \(1/4\)-separated subset \(N\) of the unit sphere.  Maximality
means that \(N\) is a \(1/4\)-net: if a unit vector were farther than \(1/4\)
from every point of \(N\), it could be added to \(N\), contradicting
maximality.

The open Euclidean balls \(B(u,1/8)\), \(u\in N\), are pairwise disjoint,
because the centers are at distance at least \(1/4\).  They are all contained
in the ball \(B(0,9/8)\), since \(\norm u=1\).  Comparing \(n\)-dimensional
Lebesgue volumes gives
\[
  |N|\left(\frac18\right)^n\vol B(0,1)
  \le
  \left(\frac98\right)^n\vol B(0,1).
\]
Thus \(|N|\le9^n\).  A complex \(N\)-dimensional Hilbert space has real
dimension \(2N\), so the bound is \(9^{2N}=81^N\).  Finally,
\(81\le e^7\), hence \(81^N\le e^{7N}\).
\end{proof}

\begin{theorem}[Concrete high-probability bounds]
There are constants \(C,C',c>0\) such that the following events have Gaussian
probability at least \(1-\exp(-cD)\):
\[
  \norm G_{\op}\le C(\sqrt D+\sqrt S),
\]
\[
  \norm{W^\pt}_{\op}\le C',
\]
and
\[
  R(G)^2\ge cDS.
\]
After normalization by \(R(G)\), these become
\[
  \norm{\Theta(G)}_{\op}\le \frac a d,\qquad
  \norm{(\Theta(G)\Theta(G)^*)^\pt}_{\op}\le \frac b{d^2}
\]
with bad probability bounded by \(\exp(-d^2/12)\) in the paper scale.
\end{theorem}

\begin{proof}
The first estimate is obtained by applying the fixed-vector Gaussian tail and
the net-cardinality lemma to a \(1/4\)-net in the sample-index sphere, and
then applying the quadratic net-lifting lemma to pass from the net to the
operator norm.  The second estimate is obtained in the same way, but the
fixed-vector tail is the quadratic-form Bernstein bound proved above, applied
to the Hermitian quadratic forms representing the Wishart and
partial-transpose Wishart observables.  The third estimate is the lower tail
for
\[
  R(G)^2=\sum_{i,\alpha}|G_{i\alpha}|^2,
\]
a sum of \(DS\) independent mean-one exponential variables; its Chernoff proof
is the lower-tail analogue of the Gamma calculation in the diagonal section.
The normalized events follow by the deterministic radial rescaling inequalities:
\[
  \norm{\Theta(G)}_{\op}=\frac{\norm G_{\op}}{R(G)},\qquad
  \norm{(\Theta\Theta^*)^\pt}_{\op}
  =\frac{\norm{(GG^*)^\pt}_{\op}}{R(G)^2}.
\]
Splitting the bad event into a mass lower-tail failure and a Gaussian/Wishart
upper-tail failure gives the final union bound.  The numerical absorption
uses \(12\log2\le d^2\) and \(e^{-D}\le e^{-d^2/12}\) when \(D=d^2\).
\end{proof}

\chapter{Surface Measure, Polar Law, and Independence}

\section{Canonical surface measure}

Let \(U(\mathcal X)\) be the unitary group of the complex Hilbert space
\(\mathcal X\).  It is compact and acts continuously on the unit sphere
\(\Sphere(\mathcal X)\).  The action is transitive.

\begin{definition}[Surface law]
The canonical surface law \(\sigma_{\mathcal X}\) is the normalized
Hausdorff/Haar surface probability measure on \(\Sphere(\mathcal X)\).
\end{definition}

\begin{lemma}[Surface-law invariance]
For every unitary \(U\in U(\mathcal X)\),
\[
  U_*\sigma_{\mathcal X}=\sigma_{\mathcal X}.
\]
\end{lemma}

\begin{proof}
The unnormalized surface measure is the restriction of the appropriate
Hausdorff measure to the sphere.  Hausdorff measure is defined from covers by
sets of small diameter.  A real isometry preserves all distances, hence
preserves the diameter of every set and sends admissible covers of \(E\) to
admissible covers of \(U(E)\) with the same total diameter cost.  Therefore it
preserves Hausdorff measure.  A unitary map is a real isometry and maps the
unit sphere onto itself.  Hence it preserves the unnormalized surface measure
of every Borel subset of the sphere.  Dividing by the total surface measure
preserves the equality, so the normalized surface law is invariant.
\end{proof}

\begin{theorem}[Uniqueness of invariant probability measure]
Let a compact topological group \(K\) act continuously and transitively on a
compact Hausdorff space \(Y\).  If \(\mu\) and \(\nu\) are \(K\)-invariant
probability measures on \(Y\), then \(\mu=\nu\).
\end{theorem}

\begin{proof}
Let \(m_K\) be normalized Haar measure on the compact group \(K\), so \(m_K\)
is a probability measure and is invariant under left multiplication.  For a
continuous function \(f\) and any \(y\in Y\), the averaged value
\[
  A_f(y)=\int_K f(k\cdot y)\,\dd m_K(k)
\]
is constant in \(y\), by transitivity and left invariance of Haar measure.
If \(\mu\) is invariant, then
\[
  \int_Y f\,\dd\mu
  =
  \int_Y\int_K f(k\cdot y)\,\dd m_K(k)\,\dd\mu(y)
  =A_f.
\]
The same holds for \(\nu\).  Thus the two measures have the same integrals
against all continuous functions.

It remains to recall why this determines the measure.  If two finite Borel
measures on a compact Hausdorff space differ, then by regularity there is a
closed set \(F\) and an open set \(O\supset F\) on which their masses differ
after an arbitrarily small enlargement.  Urysohn's lemma gives a continuous
function \(0\le h\le1\) with \(h=1\) on \(F\) and \(h=0\) outside \(O\).  The
integrals of \(h\) then separate the two measures, contradicting equality of
all continuous integrals.  Therefore \(\mu=\nu\).
\end{proof}

\section{Gaussian direction law}

\begin{lemma}[The Gaussian direction law is a probability on the sphere]
\[
  \nu_{\mathrm{dir}}(\Sphere(\mathcal X))=1.
\]
\end{lemma}

\begin{proof}
The only possible failure of \(\Theta(G)\in\Sphere(\mathcal X)\) occurs at
\(G=0\), depending on the arbitrary value chosen there.  The standard Gaussian
law has density
\[
  \pi^{-DS}\exp(-\norm G_{\HS}^2)
\]
with respect to Lebesgue measure on the real vector space underlying
\(\mathcal X\).  A singleton has Lebesgue measure \(0\), hence
\[
  \gamma(\{0\})=0.
\]
For every \(G\ne0\), \(\Theta(G)=G/\norm G_{\HS}\) has Hilbert--Schmidt norm
\(1\).  Therefore
\[
  \gamma\{G:\Theta(G)\in\Sphere(\mathcal X)\}
  \ge \gamma(\mathcal X\setminus\{0\})=1.
\]
The reverse inequality is automatic, so the direction law has total mass
\(1\) on the sphere.
\end{proof}

\begin{lemma}[Gaussian direction invariance]
For every unitary \(U\in U(\mathcal X)\),
\[
  U_*\nu_{\mathrm{dir}}=\nu_{\mathrm{dir}}.
\]
\end{lemma}

\begin{proof}
The standard Gaussian law is unitary-invariant:
\[
  U_*\gamma=\gamma.
\]
This follows directly from the radial density
\(\pi^{-DS}\exp(-\norm G_{\HS}^2)\), since unitary maps preserve
\(\norm G_{\HS}\).
Moreover, for \(G\ne0\),
\[
  \Theta(UG)=\frac{UG}{\norm{UG}}=U\frac G{\norm G}=U\Theta(G).
\]
The equality may fail only at \(G=0\), a null set.  Hence
\[
  U_*\Theta_*\gamma=\Theta_*U_*\gamma=\Theta_*\gamma.
\]
\end{proof}

\begin{theorem}[Gaussian polar law]
\[
  \Theta_*\gamma=\sigma_{\mathcal X}.
\]
\end{theorem}

\begin{proof}
Both measures are probability measures on the same sphere.  Both are invariant
under the same compact transitive unitary action.  The uniqueness theorem for
invariant probability measures gives equality.
\end{proof}

\section{Radial slices and independence}

For a Borel set \(A\subset\R_+\), define the radial slice
\[
  \gamma_A(B)=\gamma\{G\in B:\ R(G)\in A\},
\]
for Borel sets \(B\subset\mathcal X\).
Its directional pushforward is
\[
  \nu_A=\Theta_*\gamma_A.
\]

\begin{lemma}[Invariant radial slices]
For every Borel \(A\subset\R_+\), the finite measure \(\nu_A\) on the sphere
is invariant under \(U(\mathcal X)\).
\end{lemma}

\begin{proof}
The set \(\{R\in A\}\) is unitary-invariant because the radius is preserved by
unitaries.  The Gaussian is unitary-invariant, and \(\Theta\) intertwines the
action.  Therefore the pushforward slice is invariant.
\end{proof}

\begin{lemma}[Radial slice is a scalar multiple of surface law]
For every Borel \(A\subset\R_+\),
\[
  \nu_A=\gamma\{R\in A\}\,\sigma_{\mathcal X}.
\]
\end{lemma}

\begin{proof}
If \(\gamma\{R\in A\}=0\), both sides are zero.  Otherwise normalize \(\nu_A\)
to a probability measure.  It is invariant under the compact transitive
unitary action, hence equals \(\sigma_{\mathcal X}\).  Undo the normalization.
\end{proof}

\begin{theorem}[Radius--direction independence]
The random variables \(R\) and \(\Theta\) are independent.  Equivalently, for
Borel \(A\subset\R_+\) and \(B\subset\Sphere(\mathcal X)\),
\[
  \gamma\{R\in A,\ \Theta\in B\}
  =
  \gamma\{R\in A\}\,\sigma_{\mathcal X}(B).
\]
\end{theorem}

\begin{proof}
The left hand side is \(\nu_A(B)\).  Apply the radial-slice scalar-multiple
lemma.
\end{proof}

\begin{corollary}[Quadratic radial independence]
\(R^2\) and \(\Theta\) are independent.
\end{corollary}

\begin{proof}
The variable \(R^2\) is a measurable function of \(R\), and independence is
stable under measurable transformations.
\end{proof}

\chapter{Expectation Normalization}

\section{Linear radial factor}

\begin{theorem}[Sample operator-norm factorization]
All expectations below are finite.  One has
\[
  \E_\gamma\norm G_{\op}
  =
  \E_\gamma R\cdot \E_{\sigma_{\mathcal X}}\norm X_{\op}.
\]
\end{theorem}

\begin{proof}
Pointwise, \(\norm G_{\op}=R(G)\norm{\Theta(G)}_{\op}\).  The two factors are
independent by radius--direction independence, and the direction law is
\(\sigma_{\mathcal X}\).
\end{proof}

\begin{corollary}[Normalized sample expectation]
If
\[
  \E_\gamma\norm G_{\op}\le \E R\cdot \frac{C_1}{d},
  \qquad \E R>0,
\]
then
\[
  \E_{\sigma_{\mathcal X}}\norm X_{\op}\le \frac{C_1}{d}.
\]
\end{corollary}

\section{Quadratic radial factor}

\begin{theorem}[Partial-transpose operator-norm factorization]
All expectations below are finite.  One has
\[
  \E_\gamma\norm{(GG^*)^\pt}_{\op}
  =
  \E_\gamma R^2\cdot
  \E_{\sigma_{\mathcal X}}\norm{(XX^*)^\pt}_{\op}.
\]
\end{theorem}

\begin{proof}
Use
\[
  (GG^*)^\pt=R(G)^2(\Theta(G)\Theta(G)^*)^\pt
\]
and independence of \(R^2\) and \(\Theta\).
\end{proof}

\begin{corollary}[Normalized partial-transpose expectation]
If
\[
  \E_\gamma\norm{(GG^*)^\pt}_{\op}
  \le \E R^2\cdot \frac{C_2}{d^2},
  \qquad \E R^2>0,
\]
then
\[
  \E_{\sigma_{\mathcal X}}\norm{(XX^*)^\pt}_{\op}\le \frac{C_2}{d^2}.
\]
\end{corollary}

\chapter{Spherical Levy and Localized Concentration}

\section{Strong global surface Levy}

\begin{theorem}[Strong spherical Levy inequality]
Let \(K>0\), let \(f:\Sphere(\mathcal X)\to\R\) be \(K\)-Lipschitz, and let
\(M_f\) be a median.  If \(n\) is the real dimension parameter of the sphere,
then for every \(t>0\),
\[
  \sigma_{\mathcal X}\{|f-M_f|\ge t\}
  \le
  2\exp\left(-\frac{n t^2}{4K^2}\right).
\]
\end{theorem}

\begin{proof}
We use the following form of spherical isoperimetry on the normalized
Euclidean sphere: if \(A\subset\Sphere(\mathcal X)\) is Borel and
\(\sigma_{\mathcal X}(A)\ge1/2\), then for every \(r>0\),
\[
  \sigma_{\mathcal X}\{x:\dist(x,A)>r\}
  \le
  \exp\left(-\frac{n r^2}{4}\right).
\]
This is the geometric statement called the spherical Levy isoperimetric
inequality; the present theorem is the direct Lipschitz-function form.

Let
\[
  A_-=\{x:f(x)\le M_f\},\qquad A_+=\{x:f(x)\ge M_f\}.
\]
Both sets have surface measure at least \(1/2\), by the definition of a
median.  If \(f(x)>M_f+t\) and \(y\in A_-\), then
\[
  t< f(x)-f(y)\le K\dist(x,y),
\]
so \(\dist(x,A_-)>t/K\).  Hence
\[
  \{f>M_f+t\}\subset \{x:\dist(x,A_-)>t/K\}.
\]
The isoperimetric inequality gives
\[
  \sigma_{\mathcal X}\{f>M_f+t\}
  \le
  \exp\left(-\frac{n t^2}{4K^2}\right).
\]
The same argument applied to \(-f\), or equivalently to \(A_+\), gives
\[
  \sigma_{\mathcal X}\{f<M_f-t\}
  \le
  \exp\left(-\frac{n t^2}{4K^2}\right).
\]
The event \(\{|f-M_f|\ge t\}\) is contained in the union of these two one-sided
events, and the union bound gives
\[
  \sigma_{\mathcal X}\{|f-M_f|\ge t\}
  \le
  2\exp\left(-\frac{n t^2}{4K^2}\right).
\]
\end{proof}

\section{The paper-scale concentration theorem}

Let
\[
  f_k(X)=\Tr\left(\left((XX^*)^\pt\right)^k\right),
  \qquad X\in\Sphere(\mathcal X).
\]
Let the good set be
\[
  \Omega_{a,b}(d)=
  \left\{X:
    \norm X_{\op}\le \frac a d,\quad
    \norm{(XX^*)^\pt}_{\op}\le \frac b{d^2}
  \right\}.
\]

\begin{theorem}[Local Lipschitz concentration for \(f_k\)]
Fix \(\lambda>0\) and assume \(D=d^2\), \(S=s_d\), and
\(s_d/d^2\to\lambda\).  Assume also the high-probability good-set bound
\[
  \sigma_{\mathcal X}(\Omega_{a,b}(d)^c)\le
  2e^{-d^2/12},
\]
and the deterministic Lipschitz estimate on \(\Omega_{a,b}(d)\)
\[
  \Lip(f_k|_{\Omega_{a,b}(d)})\le Ck\,d^{-(2k-2)}.
\]
Then there are constants \(C',c'>0\) such that
\[
  \Prob_{\sigma_{\mathcal X}}\left\{
    |f_k-\E f_k|\ge \eps d^{-(2k-2)}
  \right\}
  \le
  C'e^{-d^2/12}
  +C'\exp\left(-c'\frac{d^4\eps^2}{k^2}\right).
\]
\end{theorem}

\begin{proof}
Apply the localized Levy theorem to \(f_k\) on the good set.  The dimension
parameter is \(n=2d^2s_d\).  Since \(s_d/d^2\to\lambda>0\), this is bounded
above and below by positive constant multiples of \(d^4\) for all large \(d\).
The Lipschitz constant is
\(Ck\,d^{-(2k-2)}\), so the Levy exponent at deviation
\(\eps d^{-(2k-2)}\) is
\[
  \frac{n\eps^2d^{-2(2k-2)}}{C^2k^2d^{-2(2k-2)}}
  \ge c\frac{d^4\eps^2}{k^2}
\]
for a constant \(c>0\) depending only on \(\lambda\) and the Lipschitz
constant.
The median-to-mean replacement is handled by the integrated-tail bound.  If
the integrated mean--median shift is at most half the target deviation, the
event centered at the mean is contained in the event centered at the median
at half scale.  Otherwise the target exponent is small enough that the right
hand side is at least \(1\), and the elementary bound that every probability
is at most \(1\) closes the branch.  Numerical constants are absorbed into
\(C'\) and \(c'\).
\end{proof}

\chapter{Final Assembly}

\begin{theorem}[Fully assembled Appendix B theorem]
Fix \(\lambda>0\).  In the balanced induced model \(D=d^2\), \(S=s_d\), with
\(s_d/d^2\to\lambda\), the following hold for all sufficiently large \(d\).
\begin{enumerate}[label=(\roman*)]
  \item The normalized sample operator norm satisfies
  \[
    \E_{\sigma_{\mathcal X}}\norm X_{\op}\le \frac{C_1}{d}.
  \]
  \item The normalized partially transposed density satisfies
  \[
    \E_{\sigma_{\mathcal X}}\norm{(XX^*)^\pt}_{\op}\le \frac{C_2}{d^2}.
  \]
  \item The two good-set events
  \[
    \norm X_{\op}\le a/d,\qquad
    \norm{(XX^*)^\pt}_{\op}\le b/d^2
  \]
  fail with probability at most \(e^{-d^2/12}\), after choosing the explicit
  constants supplied by the high-probability theorem.
  \item For each fixed \(k\), the trace-power statistic
  \[
    f_k(X)=\Tr\left(\left((XX^*)^\pt\right)^k\right)
  \]
  satisfies the localized spherical concentration estimate at the natural
  scale \(d^{-(2k-2)}\).
\end{enumerate}
\end{theorem}

\begin{proof}
The expectation bounds are the two radial-normalization corollaries, with the
Gaussian and Wishart expectation estimates supplied by the diagonal and
off-diagonal moment pipelines.  The high-probability estimates are the
concrete net, Bernstein, Wishart, and Gaussian-mass bounds transported through
the normalization identities.  The concentration estimate is the localized Levy
theorem applied to the deterministic Lipschitz scale on the intersection of
the two good sets.  The constants are those fixed in the formal assembly:
good-set exponent \(1/12\), front constant \(4\), and localized exponent
constant \(c_{\mathrm{dim}}/(64C^2)\).
\end{proof}

\appendix

\chapter{Concordance with the Formal Development}

This appendix is the only place where formal filenames are mentioned.  The
main body above is pure mathematics; this table records the promised
one-to-one correspondence between mathematical chapters and formal modules.

\begin{longtable}{>{\raggedright\arraybackslash}p{0.38\textwidth}
                  >{\raggedright\arraybackslash}p{0.54\textwidth}}
\toprule
Formal module & Mathematical content translated here\\
\midrule
\endhead
\leanfile{RandomMatrixModel},
\leanfile{GaussianModel},
\leanfile{PartialTranspose} &
Finite index sets, sample matrices, bipartite matrices, Gaussian coordinate
space, Wishart matrices, density matrices, partial transpose, and basic
dimension identities.\\
\leanfile{AppendixBConcreteModel} &
Concrete specialization \(P=Q=\mathrm{Fin}(d)\), \(\Sigma=\mathrm{Fin}(s)\),
balanced regime \(s/d^2\to\lambda\), positivity and nonemptiness facts.\\
\leanfile{ProbabilityTools},
\leanfile{AppendixB} &
Local Lipschitz predicates, McShane extension, clipping, medians, localized
Levy reduction, median-to-mean conversion, tail integration, numerical
absorption, and final paper-scale concentration wrappers.\\
\leanfile{ComplexGaussianWick},
\leanfile{TraceWickExpansion},
\leanfile{TraceWickProductExpansion} &
Complex Wick formula, monomial encoding, trace expansion as sums over closed
walks, and product expansion for partially transposed Wishart entries.\\
\leanfile{WickCountingCore},
\leanfile{WickCountingBounds},
\leanfile{AubrunMomentSpine} &
Surviving Wick pairings, equivalence relations, class counts, weighted rank,
defect, permutation fibers, and the exact rewrite of the Wick moment into
surviving contraction sums.\\
\leanfile{AppendixBAubrunProposition71},
\leanfile{AppendixBAubrunCombinatorics},
\leanfile{AppendixBAubrunGraduate} &
Aubrun tilde words, lists (3) and (7), matching conditions, nonrepetition,
innovations, compatibility, Lemma 3.6, Lemma 7.4, Lemma 7.5, the weighted
profile-count bound, and the polynomial \(P,Q\) bounds used in Propositions
7.1 and 7.3.\\
\leanfile{AppendixBAubrunMomentInput} &
Off-diagonal linear map, Hermitian structure, trace-moment power bounds,
even-moment operator-norm conversion, and off-diagonal expectation envelope.\\
\leanfile{AppendixBDiagonalGamma} &
Diagonal gamma averages, gamma maximum tails, integration of the maximum
tail, and the diagonal expectation contribution.\\
\leanfile{AppendixBPipeline},
\leanfile{AppendixBGaussianIntegrability} &
Diagonal plus off-diagonal recombination, integrability of all Gaussian
observables, and the pipeline from moment estimates to the final expectation
bound.\\
\leanfile{HighDimensionalProbability},
\leanfile{HighProbabilityBounds},
\leanfile{AppendixBWishartBridge} &
Quadratic-form Bernstein inequalities, nets, net lifting, sample operator
norm tails, Wishart-Gamma tails, Gaussian mass lower tail, normalized events,
and paper-scale high-probability bounds.\\
\leanfile{AppendixBRadialSpherical} &
Gaussian radius, squared radius, direction, measurability, no-atom facts,
direction law supported on the sphere, and abstract radial--spherical
factorization under independence.\\
\leanfile{AppendixBSurfaceMeasure} &
Unitary action on the sphere, compactness, transitivity, invariant probability
measure uniqueness, surface measure, Haar/hausdorff invariance, and the
surface probability law.\\
\leanfile{AppendixBLevyPolarBridge},
\leanfile{AppendixBPolarRadial},
\leanfile{AppendixBSphericalLevy},
\leanfile{AppendixBSphericalConcentration} &
Transport from flattened coordinates to sample matrices, equality of Gaussian
direction law and surface law, polar law, radial slice factorization,
radius--direction independence, strong global surface Levy, and localized
Levy for the exact good set.\\
\leanfile{AppendixBNormalizedExpectations},
\leanfile{AppendixBConcreteBridge} &
Concrete expectation estimates, positivity of radial factors, Gaussian and
quadratic radial factorization, normalized bounds, and concrete probability
estimates.\\
\leanfile{AppendixBFinal} &
Final no-structure assembly: normalized expectations, high-probability
events, and localized concentration with explicit constants.\\
\bottomrule
\end{longtable}

\chapter{Checklist of Closed Mathematical Inputs}

The following list records the formerly delicate points that are now theorems
in the mathematical development above.

\begin{enumerate}[label=\textbf{C\arabic*.}]
  \item The Gaussian direction law is a probability measure on the unit sphere.
  \item The canonical surface law is invariant under the same compact
  transitive unitary action.
  \item The Gaussian direction law is invariant under that action.
  \item Compact transitive uniqueness identifies the two laws:
  \(\Theta_*\gamma=\sigma_{\mathcal X}\).
  \item Radial slices are invariant and hence scalar multiples of the surface
  law.
  \item Radius and direction are independent; so are \(R^2\) and direction.
  \item Gaussian expectation bounds divide by the positive radial mean.
  \item Quadratic Wishart-Gamma expectation bounds divide by the positive
  squared radial mean.
  \item Strong global surface Levy is available for the canonical surface law.
  \item Localized Levy transfers global concentration to the good set.
  \item The two operator-norm good sets have paper-scale high probability.
  \item Aubrun's profile-count estimate gives the polynomial bound for the
  weighted Wick profiles.
  \item The off-diagonal moment pipeline gives the required expectation
  envelope.
  \item The final Appendix B assembly has no unstated structural assumptions.
\end{enumerate}

\chapter*{Bibliographic Orientation}

The proofs in the body of this document are written out directly.  The list
below is included to identify the mathematical sources and historical
provenance of the tools used: Aubrun's combinatorial organization of the
partial-transpose moment method, the classical Wick formula, the usual matrix
norm background, Haar and surface-measure language, Gamma and quadratic-form
tail inequalities, and high-dimensional net estimates.  None of the proof
steps above is intended to require opening one of these references.

\begin{thebibliography}{99}
\bibitem{aubrun}
G.~Aubrun.
\newblock Partial transposition of random states and non-centered semicircular
distributions.
\newblock \emph{Random Matrices: Theory and Applications}, 1(2), 2012.

\bibitem{aubrun-szarek}
G.~Aubrun and S.~Szarek.
\newblock \emph{Alice and Bob Meet Banach}.
\newblock American Mathematical Society, 2017.

\bibitem{bhatia}
R.~Bhatia.
\newblock \emph{Matrix Analysis}.
\newblock Graduate Texts in Mathematics 169, Springer, 1997.

\bibitem{folland-real}
G.~B. Folland.
\newblock \emph{Real Analysis: Modern Techniques and Their Applications}.
\newblock Second edition, Wiley, 1999.

\bibitem{folland-harmonic}
G.~B. Folland.
\newblock \emph{A Course in Abstract Harmonic Analysis}.
\newblock Second edition, CRC Press, 2016.

\bibitem{hsu-kakade-zhang}
D.~Hsu, S.~M. Kakade, and T.~Zhang.
\newblock A tail inequality for quadratic forms of subgaussian random vectors.
\newblock \emph{Electronic Communications in Probability}, 17, no.~52,
1--6, 2012.

\bibitem{isserlis}
L.~Isserlis.
\newblock On a formula for the product-moment coefficient of any order of a
normal frequency distribution in any number of variables.
\newblock \emph{Biometrika}, 12(1/2):134--139, 1918.

\bibitem{laurent-massart}
B.~Laurent and P.~Massart.
\newblock Adaptive estimation of a quadratic functional by model selection.
\newblock \emph{Annals of Statistics}, 28(5):1302--1338, 2000.

\bibitem{mattila}
P.~Mattila.
\newblock \emph{Geometry of Sets and Measures in Euclidean Spaces}.
\newblock Cambridge Studies in Advanced Mathematics 44, Cambridge University
Press, 1995.

\bibitem{mcshane}
E.~J. McShane.
\newblock Extension of range of functions.
\newblock \emph{Bulletin of the American Mathematical Society},
40(12):837--842, 1934.

\bibitem{vershynin}
R.~Vershynin.
\newblock \emph{High-Dimensional Probability: An Introduction with
Applications in Data Science}.
\newblock Cambridge University Press, 2018.
\end{thebibliography}

\end{document}
